\chapter{Vocabulario}

%pagina 715 de Arab Music Theory in the modern period

La música árabe, y en especial en la improvisación, hace uso de frases típicas que forman parte de la tradición. Estas frases se asocian en ocasiones a un determinado \maqamb, pero ello no significa que no puedan utilizarse en otros ellos, sin más que transponer la estructura de la frase.

En este capítulo veremos una colección de frases que pueden considerarse <<coloquiales>>, y cuyo conocimiento es es una herramienta fundamental para lograr que una improvisación suene natural y auténtica, dentro de lo que dicta la tradición.

Por sí solo, el contenido de este capítulo no sirve para elaborar un vocabulario completo y, más importante aún, saber cómo y cuando emplearlo, de la misma manera que un libro que recoja la jerga de un idioma no basta para que un hablante sepa cómo utilizar esta. Es necesaria la exposición a esa jerga en su contexto real, verla puesta en práctica, para inferir a partir de ello su utilización, de la misma manera que es necesario escuchar las frases de este capítulo dentro de sus correspondientes improvisaciones para después ser capaz de incorporarlas en las propias de manera natural. No obstante, una recopilación así de patrones melódicos a la manera de un diccionario resulta útil para disponer de un vocabulario básico y poder así entender mejor esas improvisaciones ajenas, haciendo más productivo el estudio de \taqasim y piezas de otros intérpretes.

Dividiremos los patrones de este capítulo en tres grupos principales: Inicios, cierres, y motivos melódicos.

\section{Inicios}

En esta sección se recogen fragmentos con los que es habitual comenzar un \taqsim, y que siven de presentación de su \maqamb. La apertura y cierre de un \taqsim son las partes que más se prestan al uso de <<formulas>> predefinidas, y donde el conocimiento de estas es más necesario.

Algunas de estas fórmulas estan ligadas a determinados \maqamaatb, mientras que otras son más genéricas y pueden emplearse en otros distintos.

***

Los patrones anteriores pueden emplearse no solo para arrancar un \taqsim, sino también para comenzar una frase dentro de este. Algunos patrones que se utilizan habitualmente en este caso, y que no aparecen como inicio del \taqsimb, son los siguientes:

% Patron descendente enérgico (tonica -> ghammaz)

\section{Cierres}

***

descenso a tonica
descenso a tonica con motivo
descenso a tónica con trémolo

descenso a tónica + octava + tónica
descenso a tónica + octava + tónica
descenso a tónica + sensible + tónica

doble final (primero menos definitivo, segundo más 'epico')



\section{Motivos melódicos}

Los motivos melódicos de esta sección no se asocian, por lo general, a un \maqam determinado, y pueden utilizarse en cualquiera de ellos, sin más que adaptar los intervalos y las notas correspondientes.

Se trata simplemente de frases o estructuras melódicas habituales, que por la frecuencia con la que aparecen pueden considerarse como parte de esa expresividad <<coloquial>>, y que son las que permiten al interprete crear una melodía que evoque cierta familiaridad en el oyente.

Éntiendase esta sección como un simple <<diccionario>> de esas expresiones, de interés para enriquecer el vocabulario a la hora de improvisar.

